%
% fundamental.tex -- Beweis des Fundamentalsatzes
%
% (c) 2021 Prof Dr Andreas Müller, OST Ostschweizer Fachhochschule
%
\documentclass[tikz]{standalone}
\usepackage{amsmath}
\usepackage{times}
\usepackage{txfonts}
\usepackage{pgfplots}
\usepackage{csvsimple}
\usetikzlibrary{arrows,intersections,math}
\definecolor{darkred}{rgb}{0.8,0,0}
\begin{document}
\input{zerocurves.tex}
\def\skala{0.889}
\def\dx{1}
\def\dy{1}
\def\zerolabel#1#2#3{
	\fill[color=white,opacity=0.7] ({#1*\dx},{#2*\dy}) circle[radius=0.22];
	\node at ({#1*\dx},{#2*\dy}) {$#3$};
}
\begin{tikzpicture}[>=latex,thick,scale=\skala]

\draw[->] (-6.8,0) -- (7.0,0) coordinate[label={$\operatorname{Re}$}];
\draw[->] (0,-6.8) -- (0,7.0) coordinate[label={right:$\operatorname{Im}$}];

\fill[color=gray,opacity=0.5] (0,0) circle[radius=3];
\draw[->] (0,0) -- (10:3);
\node at (10:1.9) [above] {$r$};
\fill (0,0) circle[radius=0.08];

\begin{scope}
\clip (-6.7,-6.7) rectangle (6.7,6.7);
\realA
\realB
\realC
\realD
\realE
\imagA
\imagB
\imagC
\imagD
\imagE
\end{scope}

\zeroA
\zeroB
\zeroC
\zeroD
\zeroE

\end{tikzpicture}
\end{document}

